%! AP Statistics — Unit 7, Lesson 7.2 — Warm-Up (Answer Key)
\documentclass[10pt]{article}
\usepackage{apstats-article}
\usepackage{apstats-key}
\newcommand{\TallMath}[1]{$\displaystyle #1\rule[-1.4em]{0pt}{3.2em}$}

\begin{document}

\pageheader{Unit 7, Lesson 7.2}{Warm-Up --- Answer Key}
\namedateperiod

\vspace{0.15in}

From Unit 6, the one-sample $z$-interval for a proportion has the form
$\hat p \pm z^* \sqrt{\hat p(1 - \hat p)/n}$.

\begin{enumerate}
  \item A random sample of 100 adults finds 64 support a new policy. Verify
        conditions and construct a 95\% CI for the true proportion.

        \begin{itemize}
          \item \textbf{Random:} \ans{Random sample stated.}
          \item \textbf{Large Counts:} $n\hat p = $ \ans{64}$ \ge 10$? \ans{Yes}
                \quad $n(1-\hat p) = $ \ans{36}$ \ge 10$? \ans{Yes}
          \item $\hat p = $ \ans{0.64} \quad $z^* = $ \ans{1.960}
          \item $95\%$ CI: $0.64 \pm 1.960\sqrt{0.64(0.36)/100} = 0.64 \pm 0.094$
                $= ($ \ans{0.546} $,$ \ans{0.734} $)$
        \end{itemize}

  \item Now suppose we want to estimate a population mean. $\bar x = 42$,
        $s = 8$, $n = 36$, $\sigma$ unknown.

        \begin{enumerate}[label=\alph*.]
          \item Why can't we use a $z$-interval?
                \ansline{A $z$-interval for a mean requires knowing $\sigma$. Since $\sigma$ is unknown, we cannot compute a valid $z$-statistic; we must substitute $s$, which changes the distribution of the standardized statistic.}

          \item What distribution should we use?
                \ansline{The $t$-distribution with $df = n - 1 = 35$ degrees of freedom. Its shape depends on $df$; it has heavier tails than the standard normal to account for the extra uncertainty from estimating $\sigma$ with $s$.}
        \end{enumerate}
\end{enumerate}

\end{document}
